\section{Conclusiones y recomendaciones}

\begin{enumerate}
  \item Se desarrolló una plataforma que reúne la administración del gimnasio
  y el seguimiento deportivo. Las pruebas permitieron revisar las tareas de
  miembros, rutinas, pagos y asistencia. La versión documentada en julio usa
  dos roles: entrenador y miembro.
  \item La encuesta ayudó a reconocer necesidades entre los participantes:
  cuatro de cinco responsables identificaron dificultades con cobros y pagos,
  nueve de diez miembros no consultan información del gimnasio de forma digital
  y todos calificaron la integración propuesta con 4 o 5.
  \item La preferencia de nueve miembros por el celular orientó el diseño de
  pantallas adaptables. Las funciones principales son membresías, pagos,
  asistencia, rutinas y progreso.
  \item El escenario financiero utiliza un cliente Starter de CRC 15\,000 y
  clientes Pro Gym de CRC 25\,000. Son precios propuestos por el equipo que
  deberán revisarse con clientes reales.
  \item La proyección incluye infraestructura, trabajo y soporte. Con seis
  clientes, los ingresos mensuales superarían los costos habituales del
  servicio. Para recuperar el trabajo inicial se necesita mantener la operación
  durante más tiempo, como muestra el análisis anual.
  \item El análisis anual proyecta CRC 2\,780\,000 de ingresos y
  CRC 596\,044,36 de resultado antes de impuestos, con un margen de 21,44\%.
  El valor invertido se recuperaría en el mes 8 si se cumplen los supuestos.
  El dinero disponible al cierre se estima en CRC 1\,666\,983,56. Ese saldo es
  mayor porque incluye el aporte inicial y porque el equipo trabajaría sin
  cobrar al comienzo; no sería una ganancia lista para repartir.
  \item El crecimiento comercial se apoyará en demostraciones,
  WhatsApp Business, recomendaciones y ayuda para configurar el sistema.
  Se propone destinar las primeras ganancias a formalización, seguridad e
  infraestructura.
  \item Antes de utilizar datos reales se prepararán una cuenta independiente
  para el gimnasio piloto, copias de seguridad y atención a consultas. La
  separación de información será necesaria para atender varios gimnasios.
  \item Pulso es un proyecto nuevo para ExpoTÉCNICA, iniciado académicamente en
  2025 y ampliado en 2026 por dos estudiantes que comparten el desarrollo técnico
  y comercial.
\end{enumerate}

El siguiente paso es realizar el piloto. Durante cuatro semanas se registrarán
el uso, las tareas completadas, el tiempo empleado, las dificultades y la
intención de continuar. Con esos resultados se podrán mejorar las funciones
y revisar cuánto cuesta atender a un gimnasio.

La demostración mostrará dos tareas: crear una rutina desde Planes y registrar
una membresía con su pago desde Facturación. Así se explicará cómo se relaciona
el entrenamiento del miembro con la administración del gimnasio.

\section{Declaración de uso de inteligencia artificial}

Durante el desarrollo se utilizó ChatGPT en menos del 10\% del proceso para
consultas técnicas, localización de errores y revisión ortográfica. No creó la
idea del proyecto, la investigación ni los datos. Derian Acuña Rojas e Ignacio
Masis Cascante revisaron y validaron el contenido presentado; la responsabilidad
y la autoría final pertenecen al equipo.
